\documentclass[10pt,journal,compsoc]{IEEEtran}
\usepackage{amsmath,amssymb}
\usepackage{booktabs}
\usepackage[colorlinks=true,linkcolor=blue,citecolor=blue]{hyperref}
\newtheorem{assumption}{Assumption}
\newtheorem{theorem}{Theorem}
\newtheorem{corollary}{Corollary}

\begin{document}

\title{Supplementary Material for ``Federated Multi-Center Deep SVDD\\with Lyapunov Convergence Guarantees\\for Anomaly Detection in IIoT Networks''}

\maketitle

\noindent This document provides extended theoretical derivations and discussions that complement the main paper.

\tableofcontents
\bigskip

%% ================================================================
\section{Signal Lag Handling in RAA Weights}
\label{sup:signal-lag}

The RAA weights $w_m^t$ depend on the one-round-lagged inflation signal $\Delta R_m^t$. Because $w_m^t$ is determined before the round-$t$ aggregation, it can be treated as a constant when taking conditional expectations. Under Assumption~4 (bounded heterogeneity), an explicit lower bound on $w_{\min}^t$ can be derived, guaranteeing the existence of the descent term in Theorem~2.


%% ================================================================
\section{Comparison with Standard FedAvg Convergence}
\label{sup:fedavg-comparison}

Theorem~2 differs structurally from standard convergence results for FedAvg under convex/non-convex settings. FedAvg convergence analysis tracks the global loss $\mathcal{L}(\mathbf{W}) = \sum_m p_m \mathcal{L}_m(\mathbf{W})$ and converges to the global average optimum. Such a result provides no control over individual site losses; the global average may decrease while individual site losses continue to rise.

Theorem~2 tracks $\max_m \mathcal{L}_m(\mathbf{W})$ and converges to a neighborhood of the minimax optimum. This stronger guarantee directly precludes any site's loss from diverging. The cost is a $w_{\min}^t$ factor in the convergence rate, which may make convergence slower than FedAvg. This is an inherent trade-off between fairness and efficiency.


%% ================================================================
\section{Extended DRO Motivation}
\label{sup:dro-motivation}

RAA's negative feedback mechanism relies on the inflation signal $\Delta R_m^t$ to adjust weights, with feedback intensity controlled by the hyperparameter $\gamma$. Theorem~2 guarantees convergence for fixed $\gamma$, but the optimal $\gamma$ depends on the magnitude of inter-site distribution shift. Larger shift produces more severe inflation and demands stronger feedback. In practice, the shift magnitude $\sigma_{\mathrm{het}}^2$ is unknown and may vary over time. If $\gamma$ is too small, feedback is insufficient to suppress inflation; if too large, weights oscillate.

A more fundamental difficulty is that even with a well-chosen $\gamma$, RAA's weight formula implicitly assumes that the inflation signal $\Delta R_m^t$ accurately reflects the degree of inter-site distributional conflict. Yet $\Delta R_m^t$ is only a one-round-lagged scalar observation and may produce misleading signals due to stochastic gradient noise. An accidental radius fluctuation at one site could be misinterpreted as systematic inflation, triggering excessive weight adjustment.

We introduce Wasserstein distributionally robust optimization (DRO) to address both layers of uncertainty. Rather than assuming a known distribution shift magnitude, we optimize aggregation weights for the worst-case distribution within an uncertainty set centered on the observed empirical distribution with a Wasserstein-distance radius. This makes RAA's weight design robust to distributional uncertainty.


%% ================================================================
\section{Theoretical Assumptions and Practical Gaps}
\label{sup:assumptions}

The theoretical analysis rests on several standard assumptions, and their applicability in real industrial deployments warrants discussion.

\textbf{Lipschitz continuity (Assumption~1).}
This requires the encoder to be Lipschitz continuous with respect to its parameters. For networks with ReLU activations and bounded weights, this holds naturally; the Lipschitz constant is upper-bounded by the product of per-layer spectral norms. During actual training, however, the spectral norms evolve dynamically. A practical remedy is to periodically estimate the encoder's empirical Lipschitz constant and incorporate it into RAA's feedback signal.

\textbf{Bounded gradient divergence (Assumption~4).}
This requires the inter-site gradient divergence to have an upper bound $G^2$. This bound is typically unknown before deployment, which is precisely why we introduce Wasserstein DRO. DRO converts the dependence on $G^2$ into a dependence on the uncertainty radius $\epsilon$, which can be set adaptively in a data-driven manner.

\textbf{Cluster freezing at $T_w$.}
The two-phase training protocol requires freezing cluster assignments at round $T_w$. If the industrial process undergoes mode drift over long-term operation, the frozen clustering structure may become invalid. Possible extensions include periodic re-clustering or online cluster adaptation mechanisms.


%% ================================================================
\section{Full Wasserstein DRO Dual Derivation}
\label{sup:dro-dual}

The minimax problem,
\begin{equation}
\min_{\mathbf{w} \in \Delta_M} \; \sup_{Q \in \mathcal{U}_\epsilon^t} \; \mathbb{E}_{Q}\bigl[\mathcal{L}(\mathbf{W}^t)\bigr],
\end{equation}
can be reformulated using the classical duality theorem for Wasserstein DRO. Introducing a dual variable $\mu \geq 0$, the dual form is
\begin{equation}
\min_{\mathbf{w} \in \Delta_M} \; \inf_{\mu \geq 0} \left\{ \mu \epsilon + \sum_{m=1}^{M} p_m \sup_{z} \bigl[ w_m z - \mu |z - \mathcal{L}_m(\mathbf{W}^t)| \bigr] \right\}.
\end{equation}

The inner $\sup_z$ admits a closed-form solution. When the loss function is linear in the weights, the dual problem simplifies further. For each site $m$, the inner optimization yields
\begin{equation}
\sup_{z} \bigl[ w_m z - \mu |z - \mathcal{L}_m(\mathbf{W}^t)| \bigr] =
\begin{cases}
w_m \mathcal{L}_m(\mathbf{W}^t) & \text{if } w_m \leq \mu, \\
+\infty & \text{if } w_m > \mu.
\end{cases}
\end{equation}

\subsection{Theoretical Choice of Uncertainty Radius}

Based on finite-sample Wasserstein DRO theory, when the number of observations is $M$ (the number of sites) and the confidence level is $1 - \beta$, a sufficient condition for
\begin{equation}
P\bigl(P_{\mathrm{true}} \in \mathcal{U}_\epsilon^t\bigr) \geq 1 - \beta
\end{equation}
is $\epsilon = O\bigl(M^{-1/\max(d_\ell, 2)} \cdot \log(1/\beta)^{1/\max(d_\ell, 2)}\bigr)$, where $d_\ell$ is the dimension of the loss value. Here $d_\ell = 1$ because each site's loss is a scalar, yielding a convergence rate of $O(M^{-1/2})$. In practice, it suffices to set $\epsilon \propto M^{-1/2}$.


\end{document}
